\documentclass[12pt, a4paper]{article}
\usepackage[utf8]{inputenc}
\usepackage[spanish]{babel}
\usepackage{amsmath, amssymb}
\usepackage{geometry}
\usepackage{graphicx}
\usepackage{booktabs}
\usepackage{tikz}
\usetikzlibrary{babel}
\usepackage{hyperref}

\geometry{top=2.5cm, bottom=2.5cm, left=3cm, right=3cm}

\title{\textbf{Análisis Heurístico y Exacto para el Problema del Viajante (TSP)}\\\large Evolución de Estrategias: V3, V4, V5 y V6}
\author{Traveler Project}
\date{\today}

\begin{document}

\maketitle

\begin{abstract}
Este documento detalla la evolución matemática y algorítmica de cuatro estrategias diseñadas para resolver el Problema del Viajante de Comercio (TSP) basadas en la expansión geométrica de polígonos. Se describe el funcionamiento interno, la complejidad computacional asintótica y el desempeño empírico de cada versión. La versión más reciente, \textbf{V6 (Smoothest Angle)}, introduce una metáfora física novedosa: modelar la ruta como un \textit{hilo tendido sobre pines en un tablero}, favoreciendo inserciones que minimizan la curvatura en cada punto de inserción.
\end{abstract}

\section{Introducción}
El enfoque de ``Triangle Insertion'' construye el recorrido del viajante expandiendo iterativamente un polígono inicial. Las versiones V3--V6 representan la evolución desde una heurística polinómica determinista (V3), pasando por una heurística de relajación geométrica con visión futura (V4), hasta culminar en un solver exacto mediante \textit{Branch \& Bound} (V5). La V6 reinterpreta el proceso de inserción bajo la \textbf{analogía del hilo y pines}, priorizando rutas geométricamente suaves sobre rutas de mínimo costo inmediato.

\section{Triangle Insertion V3: Convex Hull + Cheapest}

\subsection{Descripción Matemática}
La V3 es una heurística constructiva fuertemente guiada por la geometría espacial.
\begin{enumerate}
    \item \textbf{Inicialización:} Computa el Casco Convexo (Convex Hull) utilizando el algoritmo de \textit{Andrew's Monotone Chain}. Extrae el triángulo base $\triangle_{base} = \{A, B, C\} \in \text{Hull}$ que maximiza el perímetro. Esto garantiza una frontera externa que envuelve el problema, evitando auto-intersecciones tempranas.
    \item \textbf{Cheapest Insertion Global:} En cada iteración, evalúa todos los nodos no visitados $u \in U$ contra todas las aristas actuales $(i, j) \in E$. Selecciona el par $(u, (i,j))$ que minimice el costo de inserción:
    \[
    \Delta C(i, j, u) = d(i, u) + d(u, j) - d(i, j)
    \]
    \item \textbf{Post-procesamiento:} Aplica operadores de búsqueda local \textit{2-Opt} y \textit{Or-Opt} para desenredar cruces residuales.
\end{enumerate}

\subsection{Complejidad}
\begin{itemize}
    \item El Casco Convexo toma $\mathcal{O}(N \log N)$.
    \item La búsqueda del Cheapest Insertion requiere evaluar $|U|$ candidatos contra $|E|$ aristas. En el peor caso, esto cuesta $\mathcal{O}(N^2)$ por iteración. A lo largo de $N$ inserciones, la complejidad constructiva es $\mathbf{\mathcal{O}(N^3)}$.
    \item La optimización 2-Opt en el peor de los casos asintóticos es $\mathcal{O}(N^2)$, aunque suele converger rápidamente.
\end{itemize}

\subsection{Desempeño en Benchmarks y Visualización de Falla}
\begin{itemize}
    \item \textbf{Ratio global:} $1.0268$ (empate por el 1er lugar entre las heurísticas polinómicas).
    \item \textbf{Trampa de 8 puntos:} Superada perfectamente. El algoritmo encontró el óptimo matemático ($84.67$).
\end{itemize}

\textbf{¿Dónde falla la V3?}
La heurística de \textit{Cheapest Insertion} pura falla en topologías de ``Cuello de Botella'' o cuando puntos lejanos fuerzan desvíos tardíos que el 2-Opt no puede desatar. A continuación se visualiza un escenario teórico donde la V3 caería en un óptimo local por priorizar una inserción barata a corto plazo.

\begin{figure}[h]
    \centering
    \begin{tikzpicture}[scale=0.8, thick]
        % Nodos
        \node[circle, fill=black, inner sep=2pt, label=left:A] (A) at (0, 2) {};
        \node[circle, fill=black, inner sep=2pt, label=right:B] (B) at (6, 2) {};
        \node[circle, fill=black, inner sep=2pt, label=below:C] (C) at (3, 0) {};
        \node[circle, fill=black, inner sep=2pt, label=above:D] (D) at (3, 4) {};
        \node[circle, fill=red, inner sep=2pt, label=above:E (Trampa)] (E) at (3, 1.5) {};

        % Aristas óptimo local (Ruta V3)
        \draw[blue, dashed, ->] (A) -- (D) -- (B) -- (E) -- (C) -- (A);
        
        % Aristas óptimo global
        \draw[green!60!black, ->] (A) -- (E) -- (B) -- (D) -- (C) -- (A);
    \end{tikzpicture}
    \caption{Escenario de Falla (V3): El nodo E es ``barato'' de insertar desde la base (A, B, C), pero insertarlo tempranamente obliga al nodo D a tomar un desvío gigantesco por el exterior (línea punteada azul). El 2-Opt no puede arreglar esto sin romper la topología base.}
\end{figure}

\section{Triangle Insertion V4: Look-Ahead Rotation}

\subsection{Descripción Matemática}
La V4 introduce el concepto de \textit{Relajación Geométrica de Frontera}. A diferencia de la V3 que asume que el polígono actual es rígido, la V4 permite permutar aristas locales antes de realizar una inserción.

\begin{enumerate}
    \item \textbf{Métrica de Atracción:} Para evaluar si una ruta propuesta $P$ es deseable, calcula el costo del camino sumado al costo de insertar al candidato más cercano:
    \[
    S(P) = \sum_{(i,j) \in P} d(i,j) + \min_{u \in U, (i,j) \in P} \Delta C(i, j, u)
    \]
    \item \textbf{Rotación Local (Sub-vecindad 3-Opt):} Para cada triplete de nodos consecutivos $(v_k, v_{k+1}, v_{k+2}) \in P$, genera sus 5 permutaciones. Si alguna permutación resulta en un polígono modificado $P'$ tal que $S(P') < S(P)$, la rotación se aplica permanentemente.
    \item \textbf{Doble Rotación Combinada (4-Opt Restringido):} Tras la rotación simple, el algoritmo escanea el polígono en busca de cruces geométricos 2D. Si detecta que la arista $i \rightarrow i+1$ se cruza con $j \rightarrow j+1$, extrae los tríos $T_1 = \{i-1, i, i+1\}$ y $T_2 = \{j-1, j, j+1\}$. Evalúa las $6 \times 6 = 36$ permutaciones combinadas de ambos sub-polígonos. Si alguna reduce la métrica $S(P')$, se aplica la doble rotación.
    \item \textbf{Inserción:} Posteriormente, procede con la lógica \textit{Cheapest Insertion} descrita en V3.
\end{enumerate}

\subsection{Complejidad}
\begin{itemize}
    \item Existen $N$ tripletes en el polígono. Evaluar $S(P)$ toma $\mathcal{O}(N^2)$.
    \item La rotación simple de tripletes cuesta $\mathcal{O}(N^3)$ por inserción.
    \item La doble rotación detecta cruces en $\mathcal{O}(N^2)$ comprobaciones. Al encontrar un cruce, evalúa 36 configuraciones conjuntas, cada una requiriendo evaluar $S(P)$ en $\mathcal{O}(N^2)$. En el peor caso (con $\mathcal{O}(N^2)$ cruces simultáneos), esto toma $\mathcal{O}(N^4)$ por paso de inserción.
    \item Sumando ambos procesos a lo largo de $N$ inserciones, la complejidad global de la V4 optimizada se mantiene en $\mathbf{\mathcal{O}(N^5)}$ en el peor caso teórico, siendo sumamente eficiente en la práctica ya que el número de intersecciones reales es usualmente bajo ($\le 2$).
\end{itemize}

\begin{itemize}
    \item \textbf{Ratio global:} $1.0268$. Obtiene excelentes resultados con cotas muy cercanas al óptimo exacto.
    \item La adición de la Doble Rotación le otorga una robustez superior para esquivar óptimos locales complejos (dobles cruces) sin sacrificar la viabilidad polinómica del algoritmo.
\end{itemize}

\textbf{Análisis del Nudo 4-Opt Anteriormente Irresoluble}
Antes de la mejora, la V4 clásica evaluaba rotaciones de \textbf{3 nodos consecutivos} y era matemáticamente incapaz de desenredar trampas de \textbf{componentes entrelazados} (que requieren romper 4 aristas simultáneamente, es decir, un 4-Opt). Con la nueva doble rotación combinada, esta trampa geométrica se resuelve exitosamente.

\begin{figure}[h]
    \centering
    \begin{tikzpicture}[scale=0.8, thick]
        % Nodos Cuadrado Izquierdo
        \node[circle, fill=black, inner sep=2pt] (L1) at (0, 0) {};
        \node[circle, fill=black, inner sep=2pt] (L2) at (0, 3) {};
        \node[circle, fill=black, inner sep=2pt] (L3) at (3, 3) {};
        \node[circle, fill=black, inner sep=2pt] (L4) at (3, 0) {};
        
        % Nodos Cuadrado Derecho (Entrelazado)
        \node[circle, fill=black, inner sep=2pt] (R1) at (2, -1) {};
        \node[circle, fill=black, inner sep=2pt] (R2) at (2, 4) {};
        \node[circle, fill=black, inner sep=2pt] (R3) at (5, 4) {};
        \node[circle, fill=black, inner sep=2pt] (R4) at (5, -1) {};

        % Ruta Óptimo Local entrelazado (Ruta V4 atrapada)
        \draw[red, dashed] (L1) -- (L2) -- (R2) -- (R3) -- (R4) -- (R1) -- (L4) -- (L3) -- (L1);
        
        % Ruta Óptima real (Dos cuadrados separados conectados)
        \draw[green!60!black] (L1) -- (L2) -- (L3) -- (L4) -- (R1) -- (R4) -- (R3) -- (R2) -- (L1);
    \end{tikzpicture}
    \caption{Escenario de Falla (V4): Nudo 4-Opt. La línea punteada roja muestra cómo la V4 conectaría las estructuras creando un cruce doble. Al intentar rotar solo 3 nodos a la vez, cualquier movimiento intermedio empeora la distancia temporalmente. El algoritmo abandona la rotación y queda atrapado.}
\end{figure}

\section{Triangle Insertion V5: Root Backtracking (Branch \& Bound)}

\subsection{Descripción Matemática}
La V5 abandona la naturaleza heurística (\textit{Greedy}) para convertirse en un \textbf{Solver Exacto}. Utiliza la topología del problema para crear un árbol de búsqueda sistemática y exhaustiva.

\begin{enumerate}
    \item \textbf{Raíz del Árbol:} Al igual que en versiones anteriores, inicializa la búsqueda empleando el triángulo del Casco Convexo de mayor perímetro. Esto reduce dramáticamente la simetría y las redundancias del árbol de permutaciones estándar.
    \item \textbf{Cota Superior (Upper Bound):} Inicializa la mejor distancia $D_{best}$ usando una ejecución ultrarrápida del \textit{Nearest Neighbor}.
    \item \textbf{Ramificación (Branch):} Por cada estado parcial (una ruta $P_t$ y un conjunto de no visitados $U_t$), genera ramas para todas las combinaciones de insertar $u \in U_t$ en la arista $(i,j) \in P_t$.
    \item \textbf{Poda (Bound):} Si el costo de $P_t$ excede a $D_{best}$, la rama entera se descarta.
\end{enumerate}

\subsection{Complejidad}
\begin{itemize}
    \item El algoritmo navega por un espacio de búsqueda combinatorio. La complejidad asintótica en el peor de los casos (sin podas) es factorial: $\mathbf{\mathcal{O}(N!)}$.
    \item Las podas por \textit{Upper Bound} disminuyen drásticamente el espacio real explorado, permitiendo resolver problemas de hasta $\sim 15$ nodos en tiempos interactivos, pero la escalabilidad sigue estando atada a la maldición de la dimensionalidad (NP-Hard).
\end{itemize}

\subsection{Desempeño en Benchmarks}
\begin{itemize}
    \item \textbf{Ratio global:} $\mathbf{1.000}$ (Óptimo absoluto garantizado).
    \item Permitió descubrir errores metodológicos en las marcas del suite de pruebas, demostrando irrefutablemente que el límite teórico del escenario aleatorio de 8 nodos no era $71.80$, sino $84.67488$.
\end{itemize}

\section{Análisis de Falla Real de Heurísticas (8 Nodos)}
Para determinar con exactitud el límite práctico de las heurísticas constructivas polinomiales (V3 y V4), implementamos un motor de búsqueda automática que analizó 500 permutaciones aleatorias de semillas de coordenadas bidimensionales de diferentes tamaños. 

El primer escenario real de discrepancia se localizó en una topología específica de **8 nodos**, donde la naturaleza heurística (\textit{Greedy}) de la V3 y V4 no logró converger al óptimo exacto.

\begin{itemize}
    \item \textbf{Ruta Subóptima (V3 y V4):} Distancia de **1222.5**. La ruta se vio forzada a conectar los extremos de la frontera en un orden no óptimo debido a la ponderación del \textit{Cheapest Insertion} temprano.
    \item \textbf{Ruta Óptima (V5 - Exacta):} Distancia de **1216.0**. El solver exacto determinó que re-orientar la arista del cuadrante inferior reduce la distancia total a cambio de sacrificar la aparente ``baratura'' local en el paso intermedio.
\end{itemize}

\begin{figure}[ht]
    \centering
    \begin{minipage}{0.48\textwidth}
        \centering
        \begin{tikzpicture}[scale=0.8, thick]
            \draw[gray!20, thin, step=1cm] (0,0) grid (8,8);
            % Nodos reales escalados (/50)
            \node[circle, fill=black, inner sep=1.5pt, label=above right:0] (N0) at (4.12, 3.42) {};
            \node[circle, fill=black, inner sep=1.5pt, label=left:1] (N1) at (1.41, 2.62) {};
            \node[circle, fill=black, inner sep=1.5pt, label=above:2] (N2) at (1.24, 3.96) {};
            \node[circle, fill=black, inner sep=1.5pt, label=right:3] (N3) at (5.45, 0.27) {};
            \node[circle, fill=black, inner sep=1.5pt, label=above:4] (N4) at (0.78, 7.26) {};
            \node[circle, fill=black, inner sep=1.5pt, label=left:5] (N5) at (0.75, 3.98) {};
            \node[circle, fill=black, inner sep=1.5pt, label=above:6] (N6) at (6.36, 5.46) {};
            \node[circle, fill=black, inner sep=1.5pt, label=right:7] (N7) at (7.50, 5.91) {};
            
            % Ruta V3/V4: 3 -> 7 -> 6 -> 4 -> 5 -> 2 -> 1 -> 0 -> 3
            \draw[red, dashed] (N3) -- (N7) -- (N6) -- (N4) -- (N5) -- (N2) -- (N1) -- (N0) -- (N3);
        \end{tikzpicture}
        \caption{Ruta Subóptima V3/V4 (Dist: 1222.5)}
    \end{minipage}
    \hfill
    \begin{minipage}{0.48\textwidth}
        \centering
        \begin{tikzpicture}[scale=0.8, thick]
            \draw[gray!20, thin, step=1cm] (0,0) grid (8,8);
            % Nodos reales escalados (/50)
            \node[circle, fill=black, inner sep=1.5pt, label=above right:0] (N0) at (4.12, 3.42) {};
            \node[circle, fill=black, inner sep=1.5pt, label=left:1] (N1) at (1.41, 2.62) {};
            \node[circle, fill=black, inner sep=1.5pt, label=above:2] (N2) at (1.24, 3.96) {};
            \node[circle, fill=black, inner sep=1.5pt, label=right:3] (N3) at (5.45, 0.27) {};
            \node[circle, fill=black, inner sep=1.5pt, label=above:4] (N4) at (0.78, 7.26) {};
            \node[circle, fill=black, inner sep=1.5pt, label=left:5] (N5) at (0.75, 3.98) {};
            \node[circle, fill=black, inner sep=1.5pt, label=above:6] (N6) at (6.36, 5.46) {};
            \node[circle, fill=black, inner sep=1.5pt, label=right:7] (N7) at (7.50, 5.91) {};
            
            % Ruta V5 Óptima: 3 -> 0 -> 6 -> 7 -> 4 -> 5 -> 2 -> 1 -> 3
            \draw[green!60!black] (N3) -- (N0) -- (N6) -- (N7) -- (N4) -- (N5) -- (N2) -- (N1) -- (N3);
        \end{tikzpicture}
        \caption{Ruta Óptima Exacta V5 (Dist: 1216.0)}
    \end{minipage}
    \caption{Comparación de rutas mapeada sobre coordenadas reales.}
\end{figure}

\section{Triangle Insertion V6: Smoothest Angle Insertion}

\subsection{Motivación: La Analogía del Hilo y los Pines}
Imagina que los nodos del mapa son \textbf{pines clavados en un tablero} y la ruta del viajante es un \textbf{hilo} que debe pasar por todos ellos. Un hilo físico real busca naturalmente la configuración de \textit{mínima tensión}: prefiere curvas suaves sobre giros bruscos.

La heurística \textit{Cheapest Insertion} de V3 y V4 equivale a un hilo que minimiza su longitud total ignorando los ángulos que forma en cada pin. Esto genera ``nudos'' angulares que un hilo real evitaría. La V6 formaliza esta intuición física con una métrica angular.

\begin{figure}[h]
    \centering
    \begin{tikzpicture}[scale=0.95, thick]
        % Panel izquierdo: Cheapest (V4)
        \begin{scope}[xshift=0cm]
            \node at (2.5, 4.8) {\small V4 --- Cheapest (giro brusco)};
            \node[circle, fill=gray!60, inner sep=3pt] (P1) at (0,2) {};
            \node[circle, fill=gray!60, inner sep=3pt] (P2) at (1,4) {};
            \node[circle, fill=gray!60, inner sep=3pt] (P3) at (3,3) {};
            \node[circle, fill=gray!60, inner sep=3pt] (P4) at (4,1) {};
            \node[circle, fill=gray!60, inner sep=3pt] (P5) at (2,0) {};
            \draw[red!70!black, ->, line width=1.2pt] (P1) -- (P2) -- (P3) -- (P5) -- (P4) -- (P1);
            \node[red!80!black, font=\small] at (2.3, 0.4) {$\theta\approx 30^\circ$};
        \end{scope}
        % Panel derecho: Smoothest (V6)
        \begin{scope}[xshift=6.5cm]
            \node at (2.5, 4.8) {\small V6 --- Smoothest Angle (suave)};
            \node[circle, fill=gray!60, inner sep=3pt] (Q1) at (0,2) {};
            \node[circle, fill=gray!60, inner sep=3pt] (Q2) at (1,4) {};
            \node[circle, fill=gray!60, inner sep=3pt] (Q3) at (3,3) {};
            \node[circle, fill=gray!60, inner sep=3pt] (Q4) at (4,1) {};
            \node[circle, fill=gray!60, inner sep=3pt] (Q5) at (2,0) {};
            \draw[green!60!black, ->, line width=1.2pt] (Q1) -- (Q2) -- (Q3) -- (Q4) -- (Q5) -- (Q1);
            \node[green!50!black, font=\small] at (3.2, 1.8) {$\theta\approx 150^\circ$};
        \end{scope}
    \end{tikzpicture}
    \caption{La V4 inserta el nodo de menor costo inmediato, generando giros bruscos ($\theta \approx 30^\circ$). La V6 prefiere inserciones con ángulos abiertos ($\theta \approx 150^\circ$), imitando el comportamiento físico de un hilo sobre pines.}
\end{figure}

\subsection{Descripción Matemática}
La V6 hereda toda la infraestructura de inicialización y rotación de la V4, reemplazando únicamente la fase de \textit{Cheapest Insertion} por la \textbf{Inserción de Ángulo más Suave}.

\begin{enumerate}
    \item \textbf{Inicialización y Rotaciones:} Idénticas a V4 (Convex Hull, rotación local de triángulos y Doble Rotación para cruces).
    \item \textbf{Smoothest Angle Insertion:} Para cada candidato $u \in U$ y cada arista $(i, j) \in P$, el ángulo de inserción $\theta_{ij}^u$ es:
    \[
    \theta_{ij}^u = \arccos\!\left(\frac{(\vec{p}_i - \vec{p}_u) \cdot (\vec{p}_j - \vec{p}_u)}{\|\vec{p}_i - \vec{p}_u\| \cdot \|\vec{p}_j - \vec{p}_u\|}\right) \in [0, \pi]
    \]
    El score combina costo e inserción angular con factor $\alpha = 2.0$:
    \[
    \text{score}(i, j, u) = \Delta C(i,j,u) \cdot \big(1 + \alpha\,(1 + \cos\theta_{ij}^u)\big)
    \]
    Cuando $\theta = \pi$ (línea recta), la penalización es nula. Cuando $\theta = 0$ (giro en U), la penalización máxima es $4\alpha \cdot \Delta C$.
    \item \textbf{Post-procesamiento (3 capas):}
    \begin{enumerate}
        \item \textit{2-Opt} (10 iteraciones): elimina cruces de largo alcance.
        \item \textit{Or-Opt} (segmentos de 1 y 2 nodos): reposiciona nodos y pares.
        \item \textbf{Node Reinsertion (exclusivo V6):} Extrae cada nodo del tour y lo re-inserta en la posición globalmente óptima. Repite hasta convergencia. Corrige errores estructurales que 2-Opt no puede deshacer.
        \item \textit{2-Opt} final (5 iteraciones) tras la reinserción.
    \end{enumerate}
\end{enumerate}

\subsection{Complejidad}
\begin{itemize}
    \item La evaluación del score angular añade coste constante por candidato: la fase constructiva es $\mathbf{\mathcal{O}(N^3)}$, igual que V3.
    \item La \textit{Node Reinsertion} evalúa $N$ nodos $\times$ $N$ posiciones por pasada, con hasta $N$ pasadas: $\mathbf{\mathcal{O}(N^3)}$ en el peor caso.
    \item Incluyendo las rotaciones heredadas de V4, la complejidad global teórica permanece en $\mathbf{\mathcal{O}(N^5)}$, con comportamiento práctico superior a V4 en la fase de refinamiento.
\end{itemize}

\subsection{Resultados Empíricos}

Se ejecutaron 320 mapas aleatorios sobre 5 tamaños de problema (semilla fija, reproducible). La comparación directa V4 vs.~V6 mostró una ventaja creciente de V6 con el tamaño:

\begin{table}[h]
    \centering
    \begin{tabular}{rrrrrrc}
        \toprule
        $N$ & Tests & V4 gana & V6 gana & Dist.\ media V4 & Dist.\ media V6 & Mejora V6 \\
        \midrule
         15 & 100 &  0 &  2 &   1552.3 &   1552.0 & $+0.022\%$ \\
         25 & 100 &  3 &  8 &   3157.7 &   3153.8 & $+0.126\%$ \\
         50 &  60 &  5 & 13 &   8872.7 &   8853.8 & $+0.214\%$ \\
         75 &  40 &  3 & 13 &  15970.7 &  15909.4 & $+0.384\%$ \\
        100 &  20 &  7 &  6 &  24445.5 &  24437.4 & $+0.033\%$ \\
        \bottomrule
    \end{tabular}
    \caption{Comparación V4 vs.~V6 en mapas aleatorios. La mejora de V6 crece con $N$ hasta 75 nodos; a 100 nodos la Node Reinsertion encuentra menos margen libre.}
\end{table}

El caso más notable fue un mapa de 25 nodos donde V6 obtuvo una solución \textbf{121.5 unidades} más corta que V4, gracias a que la métrica angular evitó un error estructural que la reinserción luego corrigió completamente.

\section{Conclusión}
El ecosistema de estrategias ofrece soluciones para distintos contextos: \textbf{V3} provee el mejor balance entre precisión ($\le 3\%$ error promedio) y complejidad $\mathcal{O}(N^3)$, ideal para grafos grandes. \textbf{V4} añade robustez geométrica ante nudos complejos mediante rotaciones de doble triángulo. \textbf{V5} actúa como el oráculo exacto para auditar cotas en mapas pequeños. \textbf{V6} cierra el ciclo integrando una filosofía de construcción físicamente motivada (mínima curvatura angular) con un post-procesamiento de reinserción global, logrando mejoras empíricas que crecen con el tamaño del problema.

\end{document}
